% Academic CV template inspired by the uploaded CV.
% Compile with: pdflatex demo_academic_cv_template.tex
% Replace the demo text with your own information.

\documentclass[11pt,letterpaper]{article}

\usepackage[letterpaper,top=0.72in,bottom=0.62in,left=0.82in,right=0.82in]{geometry}
\usepackage[T1]{fontenc}
\usepackage{lmodern}
\usepackage{microtype}
\usepackage{hyperref}
\usepackage{fancyhdr}
\usepackage{xcolor}
\usepackage{enumitem}
\usepackage{etoolbox}

% ---------- Links ----------
\hypersetup{
  colorlinks=true,
  urlcolor=blue,
  linkcolor=blue,
  citecolor=blue,
  pdfauthor={Your Name},
  pdftitle={Your Name: Curriculum Vitae}
}

% ---------- Global typography ----------
\setlength{\parindent}{0pt}
\setlength{\parskip}{0pt}
\setlength{\emergencystretch}{2em}
\urlstyle{same}

% ---------- Repeating header on pages after the first ----------
\newcommand{\cvname}{Jannatun Naim}
\pagestyle{fancy}
\fancyhf{}
\renewcommand{\headrulewidth}{0pt}
\fancyhead[L]{\small \cvname}
\fancyhead[R]{\small \thepage}
\setlength{\headheight}{14pt}
\fancypagestyle{firstpage}{%
  \fancyhf{}
  \renewcommand{\headrulewidth}{0pt}
}

\newcommand{\researchproject}[3]{%
  \hspace*{1.55em}\parbox[t]{0.955\linewidth}{%
    \textbf{#1}\par
    \vspace{0.08em}
    {\small\itshape\color{gray} #2}\par
    \vspace{0.18em}
    {\small #3}
  }\par
  \vspace{0.75em}%
}
% ---------- Reusable formatting commands ----------
\newcommand{\cvsection}[1]{%
  \vspace{1.15em}%
  {\Large #1}\par
  \vspace{0.55em}%
}

\newcommand{\cvline}[1]{%
  \hspace*{1.55em}\parbox[t]{0.955\linewidth}{#1}\par
  \vspace{0.34em}%
}

\newcommand{\cventry}[2]{%
  \hspace*{1.55em}\parbox[t]{0.955\linewidth}{\textbf{#1} #2}\par
  \vspace{0.36em}%
}

\newcounter{publication}
\newcommand{\publication}[3]{%
  \stepcounter{publication}%
  \hspace*{1.55em}\parbox[t]{0.955\linewidth}{%
    \leavevmode\llap{\thepublication.\hspace{0.4em}}\textit{#1}\par
    #2 #3
  }\par
  \vspace{0.58em}%
}

% #1 = position title
% #2 = institution or department
% #3 = date
% #4 = description
% #1 = job title
% #2 = institution or department
% #3 = date
% #4 = description
\newcommand{\job}[4]{%
  \par
  \noindent\hspace*{1em}%
  \begin{minipage}{\dimexpr\linewidth-1em\relax}
    \textbf{#1}, #2. #3.\par
    \if\relax\detokenize{#4}\relax
    \else
      #4\par
    \fi
  \end{minipage}
  \vspace{0.65em}
}

\newcommand{\talk}[2]{%
  \hspace*{1.55em}\parbox[t]{0.955\linewidth}{%
    \textbf{#1}\par
    #2
  }\par
  \vspace{0.46em}%
}

% ---------- Document ----------
\begin{document}
\thispagestyle{firstpage}

% ---------- Name and contact block ----------
{\centering\Huge \cvname\par}
\vspace{0.45em}
{\centering\small University of Colorado Boulder\enspace\textbullet\enspace Department of Computer Science\enspace\textbullet\enspace Institute of Cognitive Science\par}
\vspace{0.5em}
{\centering\footnotesize
\href{mailto:jannatun.naim@colorado.edu}{jannatun.naim@colorado.edu}\quad\textbullet\quad
\href{https://scholar.google.com/citations?user=SVmryi8AAAAJ&hl=en}{Google Scholar}\quad\textbullet\quad
\href{https://www.colorado.edu/research/ai-institute/jannatun-naim}{iSAT Profile}\quad\textbullet\quad
\href{https://colorado.edu/ics/jannatun-naim}{ICS Profile}\par}

% ---------- Research Focus ----------

%\cvsection{Research Focus}

%My current research develops natural language processing and multimodal AI methods for understanding teaching and learning from classroom interactions. I study how teacher–student discourse, collaborative student dialogue, and instructional artifacts can be computationally modeled to identify meaningful instructional practices and support educators with scalable feedback. This includes feedback that helps teachers better understand and improve their instructional practices, as well as feedback that supports them in facilitating productive classroom discussions and engaging students in meaningful learning interactions.

%\cvsection{Research Interests}

%Natural Language Processing, AI for Education, Educational Data Mining, Multimodal Classroom Discourse Analysis



% ---------- Education ----------
\cvsection{Education}

\cvline{\textbf{Ph.D. in Computer Science}, University of Colorado Boulder, 2024--present. Supervised by Prof.\ Tamara Sumner.}

\cvline{\emph{Simultaneous with Combined Ph.D., Institute of Cognitive Science.}}

%\cvline{\textbf{M.S. in Computer Science}, University of Colorado Boulder,  expected December 2026.}

\cvline{\textbf{B.S. in Engineering, Computer Science \& Engineering}, University of Chittagong, Bangladesh, 2023.}


% ---------- Current & Ongoing Research Projects ----------
\cvsection{Current \& Ongoing Research Projects}

\researchproject
  {Belonging-Centered Instruction (BCI)}
  {Research focus within the TAPS project}
  {\begin{itemize}[leftmargin=1.3em,topsep=0pt,itemsep=0.2em,parsep=0pt]
    \item Led the NLP component of the project, developing a multi-perspective framework that integrates pedagogical talk moves, dialogue acts, and discourse relations to analyze mathematics teaching and tutoring dialogue (EDM 2025). Identified instructional patterns and the pedagogical roles of utterances outside existing talk-move categories to inform actionable feedback for educators.
    \item Developed computational methods to identify and model instructional practices that support students' sense of belonging and safety (BCI), analyzing teacher--student discourse at both the utterance and session level.
    \item Developed methods to provide scalable feedback, using BCI labels, talk-move labels, and action indicators to identify strong instructional exemplars that help educators understand and improve their practices in classroom and tutoring interactions.
  \end{itemize}}

\researchproject
  {Connection Finder}
  {Research focus within the NSF Institute for Student-AI Teaming (iSAT)}
  {\begin{itemize}[leftmargin=1.3em,topsep=0pt,itemsep=0.2em,parsep=0pt]
    \item Co-led the NLP portion of the project, developing computational methods to analyze collaborative student dialogue and connect student contributions to instructional learning goals.
    \item The project incorporates dependency dialogue analysis developed by co-authors to identify relationships between student utterances and recover the conversational context needed to interpret brief or incomplete contributions.
    \item Led the development of contextual paraphrasing informed by dialogue history, curriculum, and lesson materials to clarify references and express students' ideas more explicitly for learning goal classification.
    \item Led the development of content filtering to focus the analysis on contributions relevant to instructional content, supporting feedback that helps educators trace progress toward learning goals and facilitate productive classroom discussions.
  \end{itemize}}

\researchproject
  {Multimodal Analysis of Classroom Instruction}
  {Area exam research direction}
  {Investigated how multimodal classroom signals, including spoken discourse, whiteboard content, written annotations, gestures, and other instructional artifacts can be jointly modeled to better understand instructional quality and student engagement. A longer-term goal of this work is to develop annotation schemas and taxonomies that capture how verbal, visual, and interactional signals relate to one another, along with computational models that can learn these cross-modal relationships. These representations and models could support downstream tasks such as identifying pedagogical moves by teachers and students, tracing how instructional actions unfold across modalities, and developing richer, context-aware analyses of teaching and learning.}
  
% ---------- Publications ----------
% =========================
% Publications
% =========================
\cvsection{Publications}

{\small
\textit{* denotes shared first authorship.}
}

\vspace{0.4em}

% ---------- Journal Articles ----------
\subsection*{Journal Articles}

\publication
  {Leveraging Fusion of Sequence Tagging Models for Toxic Spans Detection}
  {\textbf{Jannatun Naim}*, Tashin Hossain*, Fareen Tasneem*, Abu Nowshed Chy, and Masaki Aono. 
  \textit{Neurocomputing}, Volume 500, pp. 688--702. Elsevier, 2022.}
  {}

% ---------- Conference Papers ----------
\subsection*{Conference Papers}

\publication
  {Towards Actionable Pedagogical Feedback: A Multi-Perspective Analysis of Mathematics Teaching and Tutoring Dialogue}
  {\textbf{Jannatun Naim}, Jie Cao, Fareen Tasneem, Jennifer Jacobs, Brent Milne, James H. Martin, and Tamara Sumner.
  In Proceedings of the 18th International Conference on Educational Data Mining (EDM), pp. 328--341, 2025.}
  {}
  
\publication
  {Harnessing Ensemble of Data Preprocessing and Hand-crafted Features for Irony Detection in Tweets}
  {Fareen Tasneem*, \textbf{Jannatun Naim}*, and Abu Nowshed Chy. 
  In Proceedings of the 23rd International Conference on Computer and Information Technology (ICCIT). IEEE, 2020.}
  {}

% ---------- Workshop and Shared-Task Papers ----------
\subsection*{Workshop and Shared-Task Papers}

\publication
  {KingsmanTrio at SemEval-2023 Task 10: Analyzing the Effectiveness of Transfer Learning Models for Explainable Online Sexism Detection}
  {\textbf{Jannatun Naim} et al. 
  In Proceedings of the 17th International Workshop on Semantic Evaluation (SemEval-2023), 2023. 
  Ranked 32nd out of 130 teams.}
  {}

\publication
  {CSECU-DSG at SemEval-2023 Task 6: Segmenting Legal Documents into Rhetorical Roles via Fine-tuned Transformer Architecture}
  {\textbf{Jannatun Naim} et al. 
  In Proceedings of the 17th International Workshop on Semantic Evaluation (SemEval-2023), 2023. 
  Ranked 22nd out of 27 teams.}
  {}

\publication
  {CSECU-DSG at SemEval-2021 Task 6: Orchestrating Multimodal Neural Architectures for Identifying Persuasion Techniques in Texts and Images}
  {\textbf{Jannatun Naim} et al. 
  In Proceedings of the 15th International Workshop on Semantic Evaluation (SemEval-2021), 2021. 
  Ranked 6th out of 15 teams.}
  {}

\publication
  {CSECU-DSG at SemEval-2021 Task 5: Leveraging Ensemble of Sequence Tagging Models for Toxic Spans Detection}
  {\textbf{Jannatun Naim} et al. 
  In Proceedings of the 15th International Workshop on Semantic Evaluation (SemEval-2021), 2021. 
  Ranked 25th out of 120 teams.}
  {}

\publication
  {CSECU-DSG at WNUT-2020 Task 2: Exploiting Ensemble of Transfer Learning and Hand-crafted Features for Identification of Informative COVID-19 English Tweets}
  {\textbf{Jannatun Naim} et al. 
  In Proceedings of the Sixth Workshop on Noisy User-generated Text (W-NUT 2020), 2020. 
  Ranked 38th out of 55 teams.}
  {}
% ---------- Employment ----------
\cvsection{Employment}

\job
  {Research Assistant}
  {University of Colorado Boulder}
  {Fall 2024 -- Summer 2026}
  {Conducted NLP research for an NSF-funded project on AI-supported instructional feedback, with a focus on belonging-centered teaching and tutoring practices.
 %The resulting research has been implemented across multiple school districts through Saga Education.
 \\
  \textit{Supervisor: Prof. Tamara Sumner}}

\job
  {Teaching Assistant}
  {University of Colorado Boulder}
  {Fall 2025, Fall 2026}
  {Teaching assistant for Introduction to Computational Thinking. Led recitation sessions, provided individual guidance on lab assignments, and graded student work. Contributed to refinements in course structure by observing student engagement and understanding and suggesting instructional adjustments.}

\job
  {Lecturer}
  {East Delta University}
  {2024}
  {School of Science, Engineering \& Technology\\
  Chittagong, Bangladesh.}


% ---------- Talks ----------
\cvsection{Talks}
\talk
  {Leveraging Fusion of Sequence Tagging Models for Toxic Spans Detection}
  {Guest lecture at Carleton College, 2025.}


% ---------- Service ----------
\cvsection{Service}

\cvline{\textbf{Mentor} for the iSAT High School Internship Program, guiding students from the Innovation Center of St.\ Vrain Valley Schools through a smaller-scale adaptation of the Connection Finder project, 2026. Provided an introductory overview of AI and large language models, including word embeddings, vector similarity, and next-token prediction. Taught prompt engineering through examples of how small changes in wording and added context influence model outputs, guiding students in refining prompts to obtain desired responses. Introduced the research task, annotation procedures, and computation of inter-annotator agreement. Guided students in annotating simulated dialogue data for on-task versus off-task contributions and alignment with lesson goals, then visualizing and interpreting patterns such as differences between speaking frequency and the number of on-topic contributions.}

\cvline{\textbf{Reviewer} at the 15 and 17th International Workshop on Semantic Evaluation (SemEval-2021, 2023).}

\cvsection{Awards and Honors}

%\cvline{Dean's and Departmental Excellence Fellowship, University of Colorado Boulder, 2024.}

%\cvline{Early Career Professional Development Fellowship, University of Colorado Boulder, 2024.}

\vspace{0.3em}

\cvline{\textit{International Accolade}}

\cvline{4th Best-Scoring Team, HerWILL Datathon, 2022.}

\vspace{0.3em}

\cvline{\textit{National Accolades}}

\cvline{3rd in Merit List, ITEE Level-2 Fundamental Information Technology Engineer (FE) Examination, 2022.}

\cvline{6th out of 105 teams, National Girls' Programming Contest, 2018.}

\cvline{10th out of 71 teams, NSU Inter-University Girls' Programming Contest, 2019.}

\cvline{20th out of 76 teams, National Girls' Programming Contest, 2019.}

\cvline{1st out of 30 participants, UCTC Grasshopper Girls' Programming Individual Contest, 2019.}

\cvline{110th out of 150 teams, National Collegiate Programming Contest (NCPC), 2020.}

\cvline{12th out of 45 teams, East Delta University Programming Contest, 2020.}

%\vfill
%\begin{center}
%  \small Last updated: \today
%\end{center}

\end{document}

% ---------- Awards ----------
\cvsection{Awards and Honors}
\cventry{Best Paper Award,}{Example Workshop at Example Conference, 2025.}
\cventry{Graduate Research Fellowship,}{Example Foundation, 2024--2027.}
\cventry{Outstanding Student Research Award,}{Department of Computer Science, Example University, 2024.}
\cventry{Dean's List,}{Example College, 2018--2020.}
\cventry{Undergraduate Research Award,}{Example College, 2019.}
